% Front matter — English thesis.
% ⚠️ Mirror of tese-pt/frontmatter/frontmatter.tex with the two abstracts SWAPPED: here the
% main `abstract` is English and `abstractotherlanguage` carries the Portuguese resumo. The
% content of both is the same text as the Portuguese tree, so a change to one has to be
% mirrored in the other.

\frontmatter
\pagestyle{plain}

\maketitlepage

%----------------------------------------------------------------------------------------
%	STATEMENT OF INTEGRITY
%----------------------------------------------------------------------------------------

\begin{soi}

\vspace{1cm}

I hereby declare having conducted this academic work with integrity.

I have not plagiarised or applied any form of undue use of information or falsification of
results along the process leading to its elaboration.

Therefore the work presented in this document is original and authored by me, having not been
previously used for any other end. The exceptions are explicitly acknowledged in the section
that addresses the ethical considerations, Section~\ref{sec:met_etica}.
Section~\ref{sec:met_ia}, which follows it, declares how Artificial Intelligence tools were
used and for what purpose.

I further declare that I have fully acknowledged the Code of Ethical Conduct of P.PORTO.

\vspace{1cm}

% The wording above is that of the official MEIA template (v2, 2026), word for word, including
% the sentence on exceptions, which applies here and therefore stays. It is the template that
% sends the declaration on the use of AI to the ethical considerations section
% (Section~\ref{sec:met_ia}) rather than to this page.
% The date follows the one on the cover (set in main.tex). It used to be \today and changed on
% its own at every compilation, which on a SIGNED declaration is worse than on the cover.
ISEP, Porto, September 2026

\end{soi}

%----------------------------------------------------------------------------------------
%	DEDICATION
%----------------------------------------------------------------------------------------
\begin{dedicatory}
To my family.
\end{dedicatory}

%----------------------------------------------------------------------------------------
%	ABSTRACT
%----------------------------------------------------------------------------------------

\begin{abstract}
When a stock price moves abruptly, retail investors need to determine whether the movement is
unusual, whether its origin lies with the company or the market, and whether comparable precedents
exist. Free applications display the percentage change, and the few that go further do not
expose the evidence behind them;
professional terminals do answer them, and are not accessible to this investor profile.

This dissertation presents InvestiGator, a financial alerting system built exclusively on free data
sources, which answers the three questions and exposes the evidence supporting each answer. The
system integrates four techniques: anomaly detection through a rolling \emph{z}-score, decomposition
of returns into market, sector and company components with beta shrinkage, precedent retrieval by
semantic similarity with measurement of the observed outcome, and a supervised triage model trained
on a purpose-built dataset.

Each component was evaluated against transparent baselines. Volatility-normalised detection proves
consistent across companies; semantic retrieval outperforms the base rate in every sector assessed;
the triage model does not outperform a volatility-only baseline, a negative result that is reported
and analysed.

By design, the system does not produce price forecasts.
\end{abstract}

\begin{abstractotherlanguage}
Quando o preço de uma ação se move de forma abrupta, o investidor particular necessita de
determinar se o movimento é invulgar, se a sua origem é a empresa ou o mercado, e se existem
precedentes. As aplicações gratuitas apresentam a variação percentual, e as poucas que vão
além disso não expõem a evidência que as sustenta; os terminais profissionais respondem, e
não estão acessíveis a este investidor.

Esta dissertação apresenta o InvestiGator, um sistema de alertas financeiros assente exclusivamente
em fontes de dados gratuitas, que responde às três questões e expõe a evidência de cada
resposta. O sistema integra quatro técnicas: deteção de anomalias por \emph{z}-score deslizante,
decomposição do retorno em mercado, setor e empresa com encolhimento do beta, recuperação de
precedentes por semelhança semântica com medição do desfecho observado, e um modelo supervisionado
de triagem treinado sobre um conjunto de dados próprio.

Cada componente foi avaliado contra linhas de base transparentes. A deteção normalizada pela
volatilidade revela-se consistente entre empresas; a recuperação semântica supera a taxa-base nos
cinco setores avaliados; o modelo de triagem não supera uma linha de base apenas com
volatilidade, resultado negativo reportado e analisado.

Por desenho, o sistema não produz previsões de preço.
\end{abstractotherlanguage}

%----------------------------------------------------------------------------------------
%	ACKNOWLEDGEMENTS
%----------------------------------------------------------------------------------------
% DRAFT IN THE AUTHOR'S VOICE, FOR HIM TO REWRITE. Acknowledgements are the one section of
% the document nobody can write for someone else: what is here is a factual base, with the
% right people and without sentimentality, and not the final wording.
%
% Two rules that governed the drafting and that should not be lost on rewriting:
%   1. THANK THOSE WHO TOOK PART, AND ONLY THOSE. People who rate the channel's alerts did
%      take part, and that feedback is reported in Section 5.6.5. What CANNOT be written is
%      thanks to participants of a controlled study, because none was run and Chapter 6
%      states so.
%   2. SPECIFICITY REPLACES EMOTION. "For the conversations that set the direction" says
%      more, and reads better, than "for all the support".
\begin{acknowledgements}
To Prof.\ Luís Gomes, supervisor, and to Rafael Silva, co-supervisor, for the conversations
that set the direction of this work, for the examples that showed me what already exists in
the field, and for the guidance on writing the document. Some of the decisions described here
were taken after one of those conversations.

To my family, for the support over these months.

To my colleagues at Sistrade, and to Sistrade, for the flexibility that made it possible to
reconcile this master's degree with work, and for the conversations about the application
while it was being built.

To those who follow the channel and rate the alerts they receive. That feedback is the only
information in this work that did not come from an automatic measurement.
\end{acknowledgements}

%----------------------------------------------------------------------------------------
%	CONTENTS
%----------------------------------------------------------------------------------------
\tableofcontents
% The lists are listed in the Contents. The List of Acronyms already was and the other three
% were not, which is the worst of the configurations: the reader learns that the lists are in
% the Contents and then does not find there the ones being looked for. The convention followed
% is that of Bruno Ribeiro's approved dissertation, whose Contents lists all of them with a
% page number. The cleardoublepage before addcontentsline ensures that the entry points to the
% page where the list BEGINS, and not to the previous one.
\cleardoublepage
\addcontentsline{toc}{chapter}{List of Figures}
\listoffigures
\cleardoublepage
\addcontentsline{toc}{chapter}{List of Tables}
\listoftables
% The List of Listings is printed again. The document has three listings, and the official
% MEIA template requires the list when they exist. The listings were added because the
% committee reads the PDF and nothing else: an asserted guarantee is an assertion, and the
% three this work makes -- the window that excludes the day assessed, the temporal separation
% enforced by a test, and the explanation that is the calculation itself -- are inspectable
% only if the code that produces them is printed here.
\cleardoublepage
\addcontentsline{toc}{chapter}{List of Listings}
\lstlistoflistings

%----------------------------------------------------------------------------------------
%	LIST OF SYMBOLS
%----------------------------------------------------------------------------------------
% Required by the checklist of the official MEIA template (v2, 2026), which marks it as
% mandatory. Only the symbols the document uses and defines are listed; the third column points
% to the equation or section where each is introduced, so that the list works as an index and
% not merely as a glossary. The List of Listings and the List of Algorithms are absent because
% the document has neither, and the template itself instructs removing those commands in that
% case.
\cleardoublepage
\addcontentsline{toc}{chapter}{List of Symbols}
\begin{symbols}{l p{0.60\textwidth} l}
$P_t$ & Closing price of session $t$ & Eq.~\ref{eq:zscore} \\
$r_t$ & Logarithmic return of day $t$, $\ln(P_t / P_{t-1})$ & Eq.~\ref{eq:zscore} \\
$\mu$, $\sigma$ & Mean and standard deviation of the returns in the window preceding the day & Eq.~\ref{eq:zscore} \\
$z_t$ & Rarity of the return of day $t$ against the preceding window & Eq.~\ref{eq:zscore} \\
\addlinespace
$r_m$, $\tilde{r}_s$ & Daily return of the market index and of the sector index & Eq.~\ref{eq:decomposicao} \\
$\beta_m$, $\beta_s$ & Historical sensitivity of the stock to the market and to the sector & Eq.~\ref{eq:decomposicao} \\
$\varepsilon$ & Part of the return explained neither by the market nor by the sector & Eq.~\ref{eq:decomposicao} \\
$\hat{\beta}$ & Beta estimate obtained by regression over the window & Eq.~\ref{eq:vasicek} \\
$\beta_{\text{ref}}$ & Reference beta towards which the estimate is shrunk & Eq.~\ref{eq:vasicek} \\
$\mathrm{SE}(\hat{\beta})$ & Standard error of the beta estimate & Eq.~\ref{eq:vasicek} \\
$\sigma^2_{\text{ref}}$ & Reference variance that sets the strength of the shrinkage & Eq.~\ref{eq:vasicek} \\
$w$ & Weight of the Vasicek shrinkage, between zero and one & Eq.~\ref{eq:vasicek} \\
\addlinespace
$\cos$ & Cosine similarity between two sentence vectors & Section~\ref{sec:met_recuperacao} \\
$h$ & Half-life of the recency decay applied to precedents & Section~\ref{sec:met_recuperacao} \\
$k$ & Number of precedents retrieved and presented & Section~\ref{sec:met_recuperacao} \\
\addlinespace
$p_{\text{bruta}}$ & Triage score before calibration & Eq.~\ref{eq:platt} \\
$p_{\text{calibrada}}$ & Triage score after Platt calibration & Eq.~\ref{eq:platt} \\
$a$, $b$ & Learned parameters of the Platt calibration & Eq.~\ref{eq:platt} \\
\addlinespace
$P$, $C$ & Precision and recall & Section~\ref{sec:av_desenho} \\
$F_1$ & Harmonic mean of precision and recall, \mbox{$2PC/(P+C)$} & Section~\ref{sec:av_desenho} \\
\end{symbols}

% nogroupskip: without it the vertical space VARIES with the initial letter -- entries sharing
% an initial end up tight together and the rest spaced apart, which reads as arbitrary spacing.
\printnoidxglossary[type=\acronymtype, title=List of Acronyms, nonumberlist,
                   nogroupskip]

\mainmatter
\pagestyle{thesis}
